Many data center proposals utilize low-priority flows to achieve various performance objectives (e.g., low latency, high utilization), relying on TCP as their preferred transport protocol. However, the suitability of TCP for such low-priority flows -- specifically, how \emph{prioritization-induced} delays in packet transmission can cause spurious timeouts and low utilization -- is relatively unexplored.  
In this paper, we conduct an empirical study to investigate the performance of TCP for low-priority flows under a wide range of realistic scenarios: use-cases (with accompanying workloads) where the performance of low-priority flows is crucial to the functioning of the overall system as well as various network loads and other network parameters. Our findings yield two key insights: 1) for several popular use-cases (e.g., network scheduling), TCP's performance for low-priority flows is within $2\times$ of a near-optimal scheme, 2) for emerging workloads that exhibit an \emph{on-off} behavior in the high priority queue (e.g., distributed ML model training), TCP's performance for low-priority flows is poor. 
Finally, we discuss and conduct preliminary evaluation to show that two simple strategies -- weighted fair queuing (WFQ) and \emph{cross-queue} congestion notification -- can substantially improve TCP's performance for low-priority flows. 